% Diagrama de bloques TikZ del ESP32-C6 / XIAO ESP32C6 — Nodo IoT Thread
% v3 — Hispanizado, rutas limpias sin cruces, bloques más espaciados

\begin{figure}[ht]
\centering
\resizebox{0.95\textwidth}{!}{%
\begin{tikzpicture}[
  % Estilos base
  blockbase/.style={rounded corners=3pt, minimum height=1.0cm, font=\scriptsize\sffamily,
    inner sep=5pt, line width=0.6pt, align=center},
  block-blue/.style={blockbase, draw=blue!65!black, fill=blue!12,
    drop shadow={shadow xshift=0.4mm, shadow yshift=-0.4mm, opacity=0.15}},
  block-purple/.style={blockbase, draw=purple!65!black, fill=purple!12,
    drop shadow={shadow xshift=0.4mm, shadow yshift=-0.4mm, opacity=0.15}},
  block-green/.style={blockbase, draw=green!50!black, fill=green!10,
    drop shadow={shadow xshift=0.4mm, shadow yshift=-0.4mm, opacity=0.15}},
  block-teal/.style={blockbase, draw=teal!65!black, fill=teal!12,
    drop shadow={shadow xshift=0.4mm, shadow yshift=-0.4mm, opacity=0.15}},
  block-orange/.style={blockbase, draw=orange!65!black, fill=orange!12,
    drop shadow={shadow xshift=0.4mm, shadow yshift=-0.4mm, opacity=0.15}},
  block-red/.style={blockbase, draw=red!60!black, fill=red!10,
    drop shadow={shadow xshift=0.4mm, shadow yshift=-0.4mm, opacity=0.15}},
  soc/.style={draw=blue!70!black, fill=blue!4, rounded corners=5pt,
    line width=1.4pt, inner sep=8pt},
  extblockbase/.style={rounded corners=3pt, minimum height=1.1cm, font=\scriptsize\sffamily,
    inner sep=5pt, line width=0.7pt, align=center,
    drop shadow={shadow xshift=0.3mm, shadow yshift=-0.3mm, opacity=0.12}},
  extblock-green/.style={extblockbase, draw=green!40!black, fill=green!5},
  extblock-purple/.style={extblockbase, draw=purple!55, fill=purple!6},
  extblock-orange/.style={extblockbase, draw=orange!55, fill=orange!6},
  extblock-brown/.style={extblockbase, draw=brown!55, fill=brown!6},
  extblock-red/.style={extblockbase, draw=red!45, fill=red!5},
  extblock-darkred/.style={extblockbase, draw=red!50, fill=red!5},
  extblock-gray/.style={extblockbase, draw=gray!55, fill=gray!6},
  conn/.style={-{Latex[length=2mm]}, line width=0.7pt, #1},
  bidi/.style={{Latex[length=2mm]}-{Latex[length=2mm]}, line width=0.7pt, #1},
]

% === Límite del SoC ===
\draw[soc] (-6.5, -2.8) rectangle (6.5, 4.8);
\node[font=\footnotesize\bfseries\sffamily, blue!70!black] at (0, 4.4) {ESP32-C6 SoC (QFN 5$\times$5\,mm)};

% === CPU ===
\node[block-blue, minimum width=4.0cm, minimum height=1.3cm] (cpu) at (0, 3.2) {};
\node[font=\scriptsize\bfseries\sffamily] at (0, 3.5) {RISC-V RV32IMAC};
\node[font=\scriptsize\sffamily] at (0, 2.95) {Núcleo único @ 160\,MHz (400\,DMIPS)};

% === Memorias ===
\node[block-purple, minimum width=1.8cm] (sram) at (-3.8, 3.2) {SRAM\\512\,KB};
\node[block-purple, minimum width=1.8cm] (rom) at (3.8, 3.2) {ROM\\320\,KB};

% === Radio 802.15.4 ===
\node[block-green, minimum width=3.2cm, minimum height=1.3cm] (radio154) at (-3.0, 1) {};
\node[font=\scriptsize\bfseries\sffamily] at (-3.0, 1.35) {Radio IEEE 802.15.4};
\node[font=\scriptsize\sffamily] at (-3.0, 0.85) {2,4\,GHz | TX: +21\,dBm};
\node[font=\scriptsize\sffamily] at (-3.0, 0.5) {RX: $-$105\,dBm};

% === Radio Wi-Fi ===
\node[block-teal, minimum width=2.8cm, minimum height=1.1cm] (radiowifi) at (3.0, 1) {};
\node[font=\scriptsize\bfseries\sffamily] at (3.0, 1.25) {Radio Wi-Fi 6};
\node[font=\scriptsize\sffamily] at (3.0, 0.85) {802.11ax 2,4\,GHz};

% === Periféricos internos — una sola fila espaciada ===
\node[block-orange, minimum width=1.6cm] (uart) at (-4.5, -0.9) {UART $\times$2};
\node[block-orange, minimum width=1.4cm] (spi)  at (-2.2, -0.9) {SPI $\times$2};
\node[block-orange, minimum width=1.2cm] (i2c)  at (0.0, -0.9) {I2C};
\node[block-orange, minimum width=1.6cm] (adc)  at (2.5, -0.9) {ADC 12 bits\\$\times$7};
\node[block-orange, minimum width=1.4cm] (gpio) at (4.8, -0.9) {GPIO\\$\times$22};

% === Seguridad ===
\node[block-red, minimum width=11.5cm, minimum height=0.8cm] (crypto) at (0, -2.0) {};
\node[font=\scriptsize\bfseries\sffamily] at (0, -1.85) {Seguridad HW: AES-256 | SHA-256 | RSA/ECC | Arranque seguro};

% === Bus interno ===
\draw[blue!25!gray, line width=2.5pt, opacity=0.4] (-5.8, 2.1) -- (5.8, 2.1);
\draw[blue!40!gray, line width=0.6pt] (-5.8, 2.1) -- (5.8, 2.1);
\node[font=\scriptsize\bfseries\sffamily, blue!40!gray, fill=blue!4, inner sep=2pt, rounded corners=1pt] at (0, 2.28) {Bus AHB/APB};
\draw[gray!35, line width=0.4pt] (cpu.south) -- (0, 2.1);
\draw[gray!35, line width=0.4pt] (sram.south) -- (-3.8, 2.1);
\draw[gray!35, line width=0.4pt] (rom.south) -- (3.8, 2.1);
\draw[gray!35, line width=0.4pt] (radio154.north) -- (-3.0, 2.1);
\draw[gray!35, line width=0.4pt] (radiowifi.north) -- (3.0, 2.1);

% === COMPONENTES EXTERNOS ===

% Antena 802.15.4 — horizontal directa
\node[extblock-green, minimum width=1.8cm] (ant154) at (-8.8, 1) {Antena PCB\\2,4\,GHz};
\draw[bidi=green!55!black] (radio154.west) -- (ant154.east);

% Memoria Flash externa — horizontal directa desde SPI
\node[extblock-purple, minimum width=1.8cm] (flash) at (-8.8, -0.9) {Flash SPI\\4\,MB};
\draw[bidi=purple!55] (spi.west) -- (flash.east);

% Transceptor RS-485 — sale por abajo de UART (ruta limpia)
\node[extblock-orange, minimum width=2cm] (rs485) at (-8.8, -2.8) {MAX485\\RS-485};
\draw[bidi=orange!55] (uart.south) -- (-4.5, -2.8) -- (rs485.east);

% Medidor
\node[extblock-brown, minimum width=2.2cm, minimum height=1.2cm] (meter) at (-12.2, -2.8) {Medidor\\Emsitech\\P2000-T\\DLMS/COSEM};
\draw[bidi=brown!55, line width=0.8pt] (rs485.west) -- (meter.east)
  node[midway, above, font=\tiny\sffamily] {9600\,bps};

% Regulador LDO — horizontal desde borde derecho
\node[extblock-red, minimum width=1.8cm] (ldo) at (8.8, -2.0) {LDO 3,3\,V\\AMS1117};
\draw[conn=red!55] (ldo.west) -- (6.7, -2.0);
\node[font=\scriptsize\sffamily, red!55!black] at (8.8, -2.85) {3,3\,V @ 800\,mA};

% Alimentación desde medidor
\node[extblock-darkred, minimum width=2cm] (pwr) at (12.2, -2.0) {Puerto aux.\\medidor\\5\,V / 100\,mA};
\draw[conn=red!65, line width=0.8pt] (pwr.west) -- (ldo.east)
  node[midway, above, font=\tiny\sffamily] {5\,V};

% Conector USB (depuración) — horizontal al GPIO
\node[extblock-gray, minimum width=1.8cm] (usb) at (8.8, -0.5) {USB-C\\(depuración)};
\draw[bidi=gray!55] (6.7, -0.5) -- (usb.west);

% === CONSUMO ===
\node[draw=red!40, fill=red!3, rounded corners=3pt, font=\scriptsize\sffamily,
  text width=3.5cm, inner sep=4pt] at (9.5, 1.5) {
  \textbf{Consumo (5\,V):}\\
  TX pico: 190\,mA\\
  RX activo: 19,3\,mA (SoC)\\
  Sistema completo: 27\,mA\\
  Reposo (\emph{idle}): 7\,mA
};

% === ETIQUETA MÓDULO XIAO ===
\node[draw=blue!35, fill=blue!3, rounded corners=3pt, font=\scriptsize\sffamily,
  text width=3cm, inner sep=4pt] at (9.5, 3.5) {
  \textbf{Módulo:} XIAO ESP32C6\\
  \textbf{Fabricante:} Seeed Studio\\
  \textbf{Dimensiones:} 21$\times$17,5\,mm\\
  \textbf{Temp.:} $-$40$^\circ$C a +85$^\circ$C
};

\end{tikzpicture}%
}% end resizebox
\caption{Diagrama de bloques del nodo IoT basado en ESP32-C6 (módulo XIAO ESP32C6, Seeed Studio). El SoC integra CPU RISC-V RV32IMAC @ 160\,MHz, radio IEEE~802.15.4 (TX +21\,dBm, RX $-$105\,dBm) para red Thread, 512\,KB SRAM y acelerador criptográfico AES-256/SHA-256/RSA. Periféricos: UART$\times$2 (consola + RS-485 vía MAX485 hacia medidor Emsitech P2000-T @ 9600\,bps), SPI (flash 4\,MB), I2C, ADC 12 bits. Alimentación: 5\,V desde puerto auxiliar del medidor, regulado a 3,3\,V con LDO. Consumo: 19,3\,mA SoC activo, 27\,mA sistema completo (incluyendo RS-485 + LDO).}
\label{fig:esp32c6-block-diagram}
\end{figure}